\documentclass[10pt,a4paper]{article}


\renewcommand{\familydefault}{\sfdefault}
\usepackage{anyfontsize}

\usepackage{xcolor}

% Marges très fines
\usepackage[margin=0.4in]{geometry}

% Personnalisation des titres de sections
\usepackage{titlesec}
\titleformat{\section}{\fontsize{13}{16}\bfseries\uppercase}{}{0em}{}[\vspace{-0.5ex}\titlerule]
\titlespacing{\section}{0pt}{1ex}{0.6ex}

% Espacement réduit pour les listes
\usepackage{enumitem}

% Gestion des liens
\usepackage[hidelinks]{hyperref}

% Suppression du numéro de page
\pagestyle{empty}

\begin{document}
\small

% --- EN-TÊTE ---
\begin{center}
    {\fontsize{28}{34}\bfseries THEO POLETTO} \\ \vspace{0.05cm}
    {\fontsize{16}{20}\selectfont Alternant Administrateur Systèmes \& Réseaux --- Recherche alternance Cybersécurité (Mastère IPSSI, Bac+5)} \\ \vspace{0.05cm}
    {\fontsize{11}{14}\itshape Fin de contrat actuel (Bac+3) : janvier 2027 | Recherche entreprise d'accueil en priorité pour la rentrée 2027} \\ \vspace{0.05cm}
    {\normalsize Nancy et alentours | Permis B | (+33) 6 95 84 99 70 | \href{mailto:theopoleto54@gmail.com}{theopoleto54@gmail.com} | \href{https://www.linkedin.com/in/theo-poletto/}{\textcolor{blue}{LinkedIn}} | \href{https://poletto-theo.vercel.app/}{\textcolor{blue}{Portfolio}}}
\end{center}

\vspace{0.1cm}

% --- PROFIL ---
\section*{Profil}
\noindent Actuellement en alternance chez Afludia (Bac+3, fin de contrat en janvier 2027), je poursuis mon parcours vers un Mastère Cybersécurité à l'IPSSI (Bac+5, rentrée 2027). Ma priorité : trouver mon entreprise d'accueil en alternance. Je reste ouvert aux différents domaines du secteur (SOC, sécurité opérationnelle, gouvernance), avec des bases en systèmes et réseaux (Active Directory, virtualisation, supervision).

% --- COMPÉTENCES TECHNIQUES ---
\section*{Compétences Techniques}
\begin{itemize}[leftmargin=0.5cm, noitemsep, topsep=0pt, partopsep=0pt]
    \item \textbf{Cybersécurité \& Supervision :} Security Onion (NSM/IDS, découverte), Zabbix, Grafana, gestion des identités Azure AD / Entra ID, segmentation réseau (VLAN).
    \item \textbf{Systèmes :} Windows Server, Linux, Active Directory, Hyper-V, Proxmox, VMware.
    \item \textbf{Réseaux :} VLAN, DHCP, DNS, NAT, Routage.
    \item \textbf{Automatisation :} PowerShell, WinRM, Scripting, Python.
    \item \textbf{Virtualisation \& Conteneurisation :} Docker, Infrastructure as Code.
    \item \textbf{Développement (secondaire) :} C\#, .NET, Blazor, SQL Server, Git.
\end{itemize}

% --- EXPÉRIENCES PROFESSIONNELLES ---
\section*{Expériences Professionnelles}

\noindent {\fontsize{12}{14}\bfseries Administrateur Systèmes \& Réseaux / Développeur (Alternance)} \hfill \textbf{2024 — Janvier 2027} \\
\textit{Afludia} \hfill \textit{France}
\begin{itemize}[leftmargin=0.5cm, noitemsep, topsep=0pt, partopsep=0pt]
    \item Déploiement et durcissement de postes Windows 11 (SSD, images PPKG, réinstallations), mon plus gros volume d'activité.
    \item Administration réseau (VLAN, DHCP/DNS, Wifi, câblage, switchs, routeurs), virtualisation Hyper-V et gestion NAS Synology.
    \item Sécurité et conformité (Patch Tuesday, suivi des CVE, hardening, gestion des accès) et supervision d'infrastructure via Zabbix/Grafana.
    \item Téléphonie VoIP (Avaya, OXE, Asterisk), scripts PowerShell d'automatisation, administration SQL Server (cluster Always On, requêtes Dapper) et sauvegardes Docker.
\end{itemize}

\vspace{0.1cm}

\noindent {\fontsize{12}{14}\bfseries Administrateur Système} \hfill \textbf{Juin 2024 — Septembre 2024} \\
\textit{ArcelorMittal} \hfill \textit{France}
\begin{itemize}[leftmargin=0.5cm, noitemsep, topsep=0pt, partopsep=0pt]
    \item Déploiement de serveurs physiques et virtuels.
    \item Administration des droits et accès via Active Directory et System Center.
    \item Supervision d'infrastructures et gestion des incidents.
\end{itemize}

% --- PROJETS TECHNIQUES ---
\section*{Projets Techniques}

\noindent {\fontsize{12}{14}\bfseries Cluster SQL Server sécurisé \& notifications Teams} \hfill \textbf{Réalisé (Afludia)}
\begin{itemize}[leftmargin=0.5cm, noitemsep, topsep=0pt, partopsep=0pt]
    \item Déploiement d'un cluster SQL Server (haute disponibilité) avec segmentation réseau via VLAN.
    \item Notifications automatisées (alertes) via webhooks Microsoft Teams. Technologies : SQL Server, VLAN, PowerShell, Teams.
\end{itemize}

\vspace{0.05cm}

\noindent {\fontsize{12}{14}\bfseries Plateforme d'administration distante} \hfill \textbf{Réalisé}
\begin{itemize}[leftmargin=0.5cm, noitemsep, topsep=0pt, partopsep=0pt]
    \item Interface web (Blazor) pour piloter des machines à distance via WinRM : scripts, déploiement, gestion multi-machines.
    \item Technologies : PowerShell, WinRM, Blazor, C\#.
\end{itemize}

\vspace{0.05cm}

\noindent {\fontsize{12}{14}\bfseries Infrastructure Docker \& Services Réseau} \hfill \textbf{Réalisé}
\begin{itemize}[leftmargin=0.5cm, noitemsep, topsep=0pt, partopsep=0pt]
    \item Services conteneurisés en environnement virtualisé : DNS (dnsmasq), gestionnaire de mots de passe (Bitwarden).
    \item Technologies : Docker, dnsmasq, Bitwarden, Linux.
\end{itemize}

\vspace{0.05cm}

\noindent {\fontsize{12}{14}\bfseries Infrastructure Multi-Systèmes (Windows/Linux)} \hfill \textbf{Réalisé}
\begin{itemize}[leftmargin=0.5cm, noitemsep, topsep=0pt, partopsep=0pt]
    \item Lab AD/DNS/DHCP (Windows Server) avec GPO de sécurité, intégration de clients Linux au domaine (realmd/sssd).
    \item RAID 1 logiciel, durcissement SSH (clé RSA + passphrase). Technologies : Active Directory, GPO, Linux, SSH.
\end{itemize}

% --- FORMATION ---
\section*{Formation}

\noindent \textbf{Mastère Cybersécurité (Bac+5)} \hfill \textbf{Rentrée 2027} \\
\textit{IPSSI} \hfill \textit{Alternance 2 ans, entreprise d'accueil à trouver en priorité}

\vspace{0.1cm}

\noindent \textbf{Bachelor Administrateur Systèmes, Réseaux \& Cybersécurité} \hfill \textbf{2024 — 2027} \\
\textit{LiveCampus} \hfill \textit{Alternance chez Afludia}

\vspace{0.1cm}

\noindent \textbf{BTS SNIR (Systèmes et Réseaux Informatiques)} \hfill \textbf{2020 — 2022} \\
\textit{Lycée Henri Loritz}

% --- CERTIFICATIONS ---
\section*{Certifications}

\noindent \textbf{CCNA 1} \hfill \textbf{Obtenu} \\
\textit{Cisco}

% --- LANGUES ---
\section*{Langues}

\noindent \textbf{Français} : Langue maternelle \hfill \textbf{Anglais} : B2

% --- CENTRES D'INTÉRÊT ---
\section*{Centres d'intérêt}

\noindent Informatique, Cinéma, Foot

\end{document}
